\documentclass[12pt,a4paper]{article}
\usepackage[margin=1in]{geometry}
\usepackage{times}
\usepackage{titlesec}
\usepackage{hyperref}
\usepackage{enumitem}
\usepackage{longtable}
\usepackage{booktabs}
\usepackage{array}
\usepackage{fancyhdr}
\usepackage{setspace}
\usepackage{float}
\usepackage{tabularx}
\usepackage{tikz}
\usetikzlibrary{shapes.geometric, arrows.meta, positioning, fit, backgrounds, calc}

% ── Edit your details here ────────────────────
\newcommand{\authorname}{Mihir Sanghvi (2427010544)}
\newcommand{\teachername}{Ms. Stuti Pandey}
% ──────────────────────────────────────────────

\hypersetup{colorlinks=true, linkcolor=black, urlcolor=black}

\pagestyle{fancy}
\fancyhf{}
\lhead{\textit{Software Requirements Specification for SEHAI}}
\rhead{\textit{Page \thepage}}
\renewcommand{\headrulewidth}{0pt}

% Section formatting to match IEEE template
\titleformat{\section}{\large\bfseries}{\thesection.}{0.5em}{}
\titleformat{\subsection}{\normalsize\bfseries}{\thesubsection}{0.5em}{}
\titleformat{\subsubsection}{\normalsize\bfseries}{\thesubsubsection}{0.5em}{}

\setstretch{1.15}

\begin{document}

% ─────────────────────────────────────────────
%  TITLE PAGE
% ─────────────────────────────────────────────
\begin{titlepage}
    \centering
    \vspace*{1.5cm}
    \rule{\linewidth}{2pt}\\[0.6cm]
    {\LARGE\bfseries SOFTWARE REQUIREMENTS\\[0.3cm] SPECIFICATION\par}
    \vspace{0.8cm}
    {\large for\par}
    \vspace{0.8cm}
    {\LARGE\bfseries SEHAI : Smart E-Health AI\par}
    \vspace{0.3cm}
    {\large\itshape AI-Powered Voice-Enabled Rural Healthcare Assistant\par}
    \vspace{1.8cm}
    {\normalsize Version 1.0\par}
    \vspace{1.2cm}
    {\normalsize Prepared by : \authorname\par}
    \vspace{0.8cm}
    {\normalsize Submitted to : \teachername\par}
    \vspace{0.8cm}
    {\normalsize April 2, 2026\par}
\end{titlepage}

% ─────────────────────────────────────────────
%  TABLE OF CONTENTS + REVISION HISTORY
% ─────────────────────────────────────────────
\pagenumbering{roman}
\tableofcontents
\newpage

\section*{Revision History}
\addcontentsline{toc}{section}{Revision History}

\begin{tabular}{|p{3.5cm}|p{2.5cm}|p{5.5cm}|p{2cm}|}
    \hline
    \textbf{Name} & \textbf{Date} & \textbf{Reason For Changes} & \textbf{Version} \\
    \hline
    \authorname & April 2, 2026 & Initial draft & 1.0 \\
    \hline
    & & & \\
    \hline
\end{tabular}

\newpage

% ─────────────────────────────────────────────
%  1. INTRODUCTION
% ─────────────────────────────────────────────
\pagenumbering{arabic}
\section{Introduction}

\subsection{Purpose}
This Software Requirements Specification (SRS) describes the complete software requirements for \textbf{SEHAI} (Smart E-Health AI), Version 1.0. SEHAI is an AI-powered, voice-enabled web application designed to strengthen rural healthcare delivery in India by digitising patient intake, enabling machine-learning-based disease triage, and managing patient referrals across the three-tier public health system (Sub-Centre, PHC, and CHC).

This document is intended to serve as the authoritative reference for development, testing, and evaluation of the system. It covers the entire product and not a subsystem.

\subsection{Document Conventions}
\begin{itemize}
    \item \textbf{Bold} text is used for terms being defined or emphasised.
    \item Requirement identifiers follow the pattern \texttt{REQ-<MODULE>-<NUMBER>} (e.g., \texttt{REQ-AUTH-01}).
    \item Priority levels are stated as \textbf{High}, \textbf{Medium}, or \textbf{Low}.
    \item Higher-level requirements do not automatically inherit priority; each requirement carries its own priority rating.
    \item \textit{TBD} is used where information is not yet available.
\end{itemize}

\subsection{Intended Audience and Reading Suggestions}
This document is intended for the following readers:

\begin{itemize}
    \item \textbf{Developers} -- Read all sections, paying particular attention to Sections 3, 4, and 6.
    \item \textbf{Project Evaluators / Faculty} -- Begin with Sections 1 and 2 for scope, then Section 4 for system features.
    \item \textbf{Testers} -- Focus on Section 4 (functional requirements) and Section 5 (non-functional requirements).
    \item \textbf{End Users (ANMs, PHC/CHC Staff)} -- Sections 2.3, 3.1, and 4 are most relevant.
\end{itemize}

It is recommended to read the document sequentially on first pass, then use the Table of Contents for targeted reference.

\subsection{Product Scope}
SEHAI is a role-based web platform for rural healthcare workers in India. It targets the existing three-tier public health infrastructure:

\begin{itemize}
    \item \textbf{Sub-Centre} -- served by Auxiliary Nurse Midwives (ANMs)
    \item \textbf{Primary Health Centre (PHC)} -- served by general practitioners
    \item \textbf{Community Health Centre (CHC)} -- served by specialist doctors
\end{itemize}

The product's goals are to:
\begin{enumerate}
    \item Reduce diagnostic delays at the sub-centre level through AI-assisted disease prediction.
    \item Enable voice-driven data entry for low-literacy healthcare workers.
    \item Provide a structured, trackable referral pathway from Sub-Centre to PHC to CHC.
    \item Maintain a centralised, secure patient record accessible across facilities.
\end{enumerate}

SEHAI does not replace clinical judgment; it serves as a decision-support and workflow-management tool.

\subsection{References}
\begin{enumerate}
    \item IEEE Std 830-1998, \textit{IEEE Recommended Practice for Software Requirements Specifications}.
    \item Karl E. Wiegers, \textit{SRS Template}, Process Impact, 1999.
    \item XGBoost Documentation, \url{https://xgboost.readthedocs.io}
    \item FastAPI Documentation, \url{https://fastapi.tiangolo.com}
    \item Supabase Documentation, \url{https://supabase.com/docs}
    \item React 18 Documentation, \url{https://reactjs.org}
    \item Ministry of Health and Family Welfare, India, \textit{Rural Health Statistics}, 2023.
\end{enumerate}

% ─────────────────────────────────────────────
%  2. OVERALL DESCRIPTION
% ─────────────────────────────────────────────
\section{Overall Description}

\subsection{Product Perspective}
SEHAI is a new, self-contained product with no predecessor system. It is designed as an intelligent overlay on India's existing three-tier rural healthcare structure, which currently relies on paper-based records and manual referrals.

The system consists of three independently deployable components:
\begin{itemize}
    \item \textbf{Frontend SPA} -- React/TypeScript application deployed on Vercel.
    \item \textbf{Backend API} -- FastAPI Python service hosting the XGBoost ML model, deployed on HuggingFace Spaces (Docker, port 7860).
    \item \textbf{Database} -- Supabase-hosted PostgreSQL with built-in authentication and Row-Level Security.
\end{itemize}

The frontend communicates with the backend via REST (JSON over HTTPS) and with Supabase via the Supabase JS client library.

\subsection{Product Functions}
The major functions SEHAI must perform are:
\begin{enumerate}
    \item \textbf{User Registration and Login} -- Role-based access for ANM, PHC doctor, and CHC specialist.
    \item \textbf{Voice-Enabled Patient Entry} -- Symptom input via microphone (Web Speech API).
    \item \textbf{AI Disease Prediction} -- XGBoost classifier over 133 standardised symptoms.
    \item \textbf{Automated Triage and Priority Assignment} -- Low / Medium / High / Critical classification.
    \item \textbf{Patient Referral Workflow} -- Sub-Centre $\rightarrow$ PHC $\rightarrow$ CHC escalation with status tracking.
    \item \textbf{Centralised Patient Records} -- Demographics, vitals, symptoms, history, predictions, referral chain.
    \item \textbf{Role-Specific Dashboards} -- Tailored statistics and views per user role.
    \item \textbf{Multilingual Support} -- English and Hindi (regional languages planned).
\end{enumerate}

\subsection{User Classes and Characteristics}

\begin{center}
\begin{tabular}{|p{2.8cm}|p{5cm}|p{2.2cm}|p{3.2cm}|}
    \hline
    \textbf{User Class} & \textbf{Description} & \textbf{Tech Level} & \textbf{Primary Functions Used} \\
    \hline
    ANM & Frontline health worker at Sub-Centre. First point of patient contact. & Low & Patient entry, voice input, referral to PHC \\
    \hline
    PHC Doctor & General practitioner managing referrals from multiple Sub-Centres. & Medium & Review referrals, update status, escalate to CHC \\
    \hline
    CHC Specialist & Senior doctor handling complex and critical cases from PHCs. & Medium--High & Manage critical cases, update outcomes \\
    \hline
    System Admin & Manages deployment, database, and user accounts. & High & All system functions \\
    \hline
\end{tabular}
\end{center}

\subsection{Operating Environment}
\begin{itemize}
    \item \textbf{Client:} Any modern browser (Chrome 90+, Firefox 88+, Edge 90+, Safari 14+) on desktop or mobile devices.
    \item \textbf{Backend Server:} Python 3.10+ runtime inside a Docker container on HuggingFace Spaces.
    \item \textbf{Database:} Supabase cloud-hosted PostgreSQL (managed service).
    \item \textbf{Network:} Requires internet connectivity. Low-bandwidth optimisation is desired for rural deployments.
    \item \textbf{Devices:} Android smartphones (Android 8+), low-end tablets, and desktop workstations.
    \item \textbf{Operating Systems:} Platform-agnostic (browser-based); no native installation required.
\end{itemize}

\subsection{Design and Implementation Constraints}
\begin{itemize}
    \item The ML model is pre-trained; retraining is out of scope for Version 1.0.
    \item The symptom vocabulary is fixed at 133 terms; free-form symptoms outside this set are not supported.
    \item Authentication must use Supabase Auth; third-party OAuth (Google, GitHub) is not required in v1.0.
    \item All Supabase tables must have Row-Level Security (RLS) enabled at all times.
    \item The backend API must support CORS for \texttt{localhost:5173} (development) and \texttt{*.vercel.app} (production).
    \item The system must handle patient health data in accordance with applicable Indian data protection norms.
    \item Frontend must be built with React 18 + TypeScript + Vite; no alternative framework is permitted.
\end{itemize}

\subsection{User Documentation}
The following documentation will be delivered with the system:
\begin{itemize}
    \item \textbf{In-app Help Page} (\texttt{/help}) -- Step-by-step guidance for each user role, accessible from all dashboards.
    \item \textbf{API Documentation} -- Auto-generated Swagger/OpenAPI docs at \texttt{/docs} on the FastAPI backend.
    \item \textbf{README} -- Developer setup and deployment instructions in the project repository.
\end{itemize}

\subsection{Assumptions and Dependencies}
\begin{itemize}
    \item Supabase service availability is assumed; there is no offline fallback for database operations.
    \item Voice input requires a device with a microphone and a browser that supports the Web Speech API (Chrome is recommended).
    \item The XGBoost model's triage predictions are assumed to be suitable for decision support, not definitive medical diagnosis.
    \item Regional language transcription quality depends on the availability and accuracy of third-party speech recognition services.
    \item ANMs are assumed to have received basic training on using a smartphone application.
    \item The project depends on continued free-tier availability of HuggingFace Spaces and Supabase.
\end{itemize}

% ─────────────────────────────────────────────
%  3. EXTERNAL INTERFACE REQUIREMENTS
% ─────────────────────────────────────────────
\section{External Interface Requirements}

\subsection{User Interfaces}
The system provides the following screens:

\begin{itemize}
    \item \textbf{Landing Page (\texttt{/}):} Marketing page with feature highlights, hero call-to-action, and links to Login and Signup. Includes a light/dark mode toggle.
    \item \textbf{Login (\texttt{/login}):} Fields for email, password, and role selection (ANM/PHC/CHC). Displays an inline error on authentication failure without revealing which field is wrong.
    \item \textbf{Signup (\texttt{/signup}):} Registration form with full name, email, password, confirm password, role, facility name, and facility location.
    \item \textbf{ANM Dashboard (\texttt{/anm-dashboard}):} Patient entry form, voice input button, AI prediction result panel, refer-to-PHC button, and health tips sidebar.
    \item \textbf{PHC Dashboard (\texttt{/phc-dashboard}):} Referral table (with patient ID, name, ANM, symptoms, priority badge, status badge), summary statistics cards, and case review panel.
    \item \textbf{CHC Dashboard (\texttt{/chc-dashboard}):} Specialist referral table (with case ID, patient, referring PHC, priority, status), summary statistics, and critical case management panel.
    \item \textbf{Patient Record (\texttt{/patient/:id}):} Full patient view -- demographics, vitals, symptoms, AI prediction, referral history, and action buttons (consult, prescribe, refer, request labs).
    \item \textbf{Help (\texttt{/help}):} Role-specific guidance and FAQ.
    \item \textbf{Not Found (\texttt{/404}):} Friendly error page with navigation back to dashboard.
\end{itemize}

All screens must be fully responsive (minimum 360px width) and support light and dark themes. Priority badges use colour coding: \textcolor{red}{\textbf{red}} for critical/high, \textcolor{orange}{\textbf{orange}} for medium, \textcolor{green}{\textbf{green}} for low.

\subsection{Hardware Interfaces}
\begin{itemize}
    \item \textbf{Microphone:} Required for voice-enabled symptom entry. Accessed via the browser's Web Speech API; no special drivers needed.
    \item \textbf{Display:} Minimum 360\,px width; optimised for 720p and above.
    \item \textbf{Camera, GPS, Bluetooth:} Not used in Version 1.0.
\end{itemize}

\subsection{Software Interfaces}
\begin{itemize}
    \item \textbf{Supabase JS Client (\texttt{@supabase/supabase-js} v2):} Used for authentication (JWT) and all database read/write operations via PostgREST.
    \item \textbf{FastAPI Backend (v0.115.0):} Consumed by the frontend via \texttt{fetch} or Axios. Provides three REST endpoints: \texttt{GET /health}, \texttt{GET /symptoms}, \texttt{POST /predict}.
    \item \textbf{Web Speech API:} Browser-native API; no external library or API key required.
    \item \textbf{XGBoost Runtime (v1.6.2):} Loaded by the FastAPI backend from \texttt{xgboost\_model.json} at startup.
    \item \textbf{Vite Dev Server (v6.0.1):} Used during development; Vercel serves the production build.
\end{itemize}

\subsection{Communications Interfaces}
\begin{itemize}
    \item All client-server communication is over \textbf{HTTPS} (TLS 1.2 or higher).
    \item Data format for all API calls is \textbf{JSON}.
    \item Allowed CORS origins on the backend: \texttt{http://localhost:5173} (dev) and \texttt{https://*.vercel.app} (prod).
    \item Supabase real-time websocket channels may be used for live referral status updates (planned, v1.1).
    \item No email or SMS notification interface is included in Version 1.0.
\end{itemize}

% ─────────────────────────────────────────────
%  4. SYSTEM FEATURES
% ─────────────────────────────────────────────
\section{System Features}

\subsection{User Authentication and Role Management}

\subsubsection{Description and Priority}
Provides secure registration, login, and role-based routing for ANMs, PHC doctors, and CHC specialists. \textbf{Priority: High.}

\subsubsection{Stimulus/Response Sequences}
\begin{itemize}
    \item \textit{User navigates to \texttt{/signup}} $\rightarrow$ System presents registration form $\rightarrow$ User submits $\rightarrow$ System creates profile; redirects to role dashboard.
    \item \textit{User navigates to \texttt{/login}} $\rightarrow$ Enters credentials $\rightarrow$ System validates with Supabase Auth $\rightarrow$ On success: redirect to role dashboard; on failure: display error.
    \item \textit{User clicks Logout} $\rightarrow$ System invalidates session token $\rightarrow$ Redirect to landing page.
    \item \textit{Unauthenticated user accesses dashboard URL} $\rightarrow$ System redirects to \texttt{/login}.
\end{itemize}

\subsubsection{Functional Requirements}
\begin{itemize}
    \item[\texttt{REQ-AUTH-01}] The system shall allow a new user to register with full name, email, password, role, facility name, and facility location.
    \item[\texttt{REQ-AUTH-02}] The system shall reject registration if the email is already in use, displaying an appropriate message.
    \item[\texttt{REQ-AUTH-03}] The system shall authenticate users via email and password using Supabase Auth.
    \item[\texttt{REQ-AUTH-04}] After successful login, the system shall redirect the user to the dashboard matching their role (\texttt{/anm-dashboard}, \texttt{/phc-dashboard}, or \texttt{/chc-dashboard}).
    \item[\texttt{REQ-AUTH-05}] The system shall prevent users from accessing dashboards outside their assigned role.
    \item[\texttt{REQ-AUTH-06}] The system shall provide a logout function that invalidates the session and clears local state.
\end{itemize}

\subsection{Voice-Enabled Patient Entry (ANM)}

\subsubsection{Description and Priority}
Allows ANMs to register a new patient and enter symptoms by voice, with AI-assisted disease prediction triggered on submission. \textbf{Priority: High.}

\subsubsection{Stimulus/Response Sequences}
\begin{itemize}
    \item \textit{ANM clicks ``Add Patient''} $\rightarrow$ System shows patient entry form.
    \item \textit{ANM clicks microphone icon} $\rightarrow$ System activates Web Speech API $\rightarrow$ ANM speaks symptoms $\rightarrow$ System transcribes and maps to standardised symptom tags.
    \item \textit{ANM submits form} $\rightarrow$ System saves patient record $\rightarrow$ Calls \texttt{POST /predict} $\rightarrow$ Displays predicted disease, confidence, description, and precautions.
\end{itemize}

\subsubsection{Functional Requirements}
\begin{itemize}
    \item[\texttt{REQ-ANM-01}] The system shall provide a form for patient name, age, gender, symptoms (text), and clinical notes.
    \item[\texttt{REQ-ANM-02}] The system shall provide a microphone button that activates the Web Speech API for hands-free symptom dictation.
    \item[\texttt{REQ-ANM-03}] The system shall map transcribed speech to the 133-symptom standardised vocabulary where possible.
    \item[\texttt{REQ-ANM-04}] On form submission, the system shall save the patient record to the \texttt{patients} table and return a generated patient ID.
    \item[\texttt{REQ-ANM-05}] The system shall call \texttt{POST /predict} with the entered symptoms and display the result (disease, confidence, description, precautions) on the dashboard.
    \item[\texttt{REQ-ANM-06}] The prediction result shall be saved to the \texttt{predictions} table linked to the patient.
    \item[\texttt{REQ-ANM-07}] The ANM dashboard shall display contextual health tips (malaria prevention, hydration, vaccination, hand hygiene) in a sidebar.
\end{itemize}

\subsection{Patient Referral to PHC}

\subsubsection{Description and Priority}
Allows ANMs to escalate patients to a PHC with priority classification and reason documentation. \textbf{Priority: High.}

\subsubsection{Stimulus/Response Sequences}
\begin{itemize}
    \item \textit{ANM clicks ``Refer to PHC''} for a patient $\rightarrow$ System shows referral form $\rightarrow$ ANM selects target PHC, priority, and enters reason $\rightarrow$ System creates referral record with status \texttt{pending}.
\end{itemize}

\subsubsection{Functional Requirements}
\begin{itemize}
    \item[\texttt{REQ-REF-01}] The system shall allow an ANM to create a referral for any patient they registered, specifying target facility, priority (low/medium/high/critical), and reason.
    \item[\texttt{REQ-REF-02}] A new referral shall be saved to the \texttt{referrals} table with status \texttt{pending} and a creation timestamp.
    \item[\texttt{REQ-REF-03}] The referred patient shall appear in the target PHC doctor's referral table immediately after creation.
\end{itemize}

\subsection{PHC Dashboard and Referral Management}

\subsubsection{Description and Priority}
Provides PHC doctors with a view of all incoming referrals, tools to review cases, update status, and escalate to CHC. \textbf{Priority: High.}

\subsubsection{Stimulus/Response Sequences}
\begin{itemize}
    \item \textit{PHC doctor logs in} $\rightarrow$ Dashboard loads referral table with all pending/active referrals for their facility.
    \item \textit{Doctor clicks a referral row} $\rightarrow$ System opens full patient record.
    \item \textit{Doctor updates referral status} $\rightarrow$ System updates \texttt{referrals} table with new status, reviewer ID, and timestamp.
    \item \textit{Doctor clicks ``Refer to CHC''} $\rightarrow$ System creates a new referral record targeting a CHC.
\end{itemize}

\subsubsection{Functional Requirements}
\begin{itemize}
    \item[\texttt{REQ-PHC-01}] The PHC dashboard shall display all referrals addressed to the doctor's facility, showing patient ID, name, referring ANM, symptoms, priority badge, and status badge.
    \item[\texttt{REQ-PHC-02}] The dashboard shall display summary statistics: total cases, pending referrals, treated cases, and cases referred to CHC.
    \item[\texttt{REQ-PHC-03}] The doctor shall be able to open a full patient record from any referral row.
    \item[\texttt{REQ-PHC-04}] The doctor shall be able to update a referral status to \texttt{reviewed}, \texttt{in\_progress}, or \texttt{completed}, with optional notes.
    \item[\texttt{REQ-PHC-05}] The doctor shall be able to escalate a patient to a CHC by creating a new referral with updated priority and clinical notes.
\end{itemize}

\subsection{CHC Dashboard and Specialist Management}

\subsubsection{Description and Priority}
Provides CHC specialists with a view of escalated cases and tools to manage critical patients. \textbf{Priority: High.}

\subsubsection{Stimulus/Response Sequences}
\begin{itemize}
    \item \textit{CHC specialist logs in} $\rightarrow$ Dashboard loads specialist referral table.
    \item \textit{Specialist updates case status} $\rightarrow$ System saves new status and timestamp.
\end{itemize}

\subsubsection{Functional Requirements}
\begin{itemize}
    \item[\texttt{REQ-CHC-01}] The CHC dashboard shall display all referrals targeted at the specialist's facility, showing case ID, patient details, referring PHC, symptoms, priority, and status.
    \item[\texttt{REQ-CHC-02}] The dashboard shall display summary statistics: total cases, critical cases, completed, and in-progress.
    \item[\texttt{REQ-CHC-03}] The specialist shall be able to update case status to \texttt{in\_progress} or \texttt{completed} and add consultation notes.
\end{itemize}

\subsection{AI Disease Prediction Engine}

\subsubsection{Description and Priority}
A FastAPI backend service that runs an XGBoost classifier over 133 standardised symptoms to return a predicted disease with confidence score and precautions. \textbf{Priority: High.}

\subsubsection{Stimulus/Response Sequences}
\begin{itemize}
    \item \textit{Frontend sends \texttt{POST /predict} with symptom array} $\rightarrow$ Backend encodes to 133-dim binary vector $\rightarrow$ XGBoost inference $\rightarrow$ Returns predicted disease, confidence, description, precautions.
    \item \textit{Frontend sends \texttt{GET /symptoms}} $\rightarrow$ Backend returns list of 133 supported symptom names.
    \item \textit{Frontend sends \texttt{GET /health}} $\rightarrow$ Backend returns service status and model load state.
\end{itemize}

\subsubsection{Functional Requirements}
\begin{itemize}
    \item[\texttt{REQ-AI-01}] The \texttt{POST /predict} endpoint shall accept an array of symptom name strings and return: predicted disease name, confidence score (0--1), disease description, and list of precautions.
    \item[\texttt{REQ-AI-02}] The endpoint shall return HTTP 400 if the symptom array is empty.
    \item[\texttt{REQ-AI-03}] The endpoint shall return HTTP 500 if the XGBoost model failed to load at startup.
    \item[\texttt{REQ-AI-04}] The \texttt{GET /symptoms} endpoint shall return the complete list of 133 supported symptom names and their count.
    \item[\texttt{REQ-AI-05}] The \texttt{GET /health} endpoint shall return service status and whether the model is loaded.
\end{itemize}

\subsection{Patient Record Management}

\subsubsection{Description and Priority}
Provides a comprehensive, role-accessible patient record view with actions. \textbf{Priority: Medium.}

\subsubsection{Stimulus/Response Sequences}
\begin{itemize}
    \item \textit{Authorised user navigates to \texttt{/patient/:id}} $\rightarrow$ System loads full patient record from Supabase.
    \item \textit{User clicks an action button} $\rightarrow$ System processes the action (update notes, create referral, etc.).
\end{itemize}

\subsubsection{Functional Requirements}
\begin{itemize}
    \item[\texttt{REQ-PAT-01}] The patient record page shall display: name, age, gender, contact, address, symptoms, clinical notes, medical history, vitals (temperature, BP, heart rate, SpO2), AI prediction results, and referral history.
    \item[\texttt{REQ-PAT-02}] Authorised users shall be able to start a consultation, add prescription notes, request lab tests, or create a CHC referral from the patient record page.
    \item[\texttt{REQ-PAT-03}] Only users with a valid session and appropriate role shall be able to access patient records (enforced via RLS).
\end{itemize}

% ─────────────────────────────────────────────
%  5. OTHER NONFUNCTIONAL REQUIREMENTS
% ─────────────────────────────────────────────
\section{Other Nonfunctional Requirements}

\subsection{Performance Requirements}
\begin{itemize}
    \item \texttt{POST /predict} shall respond within \textbf{2 seconds} under normal load conditions.
    \item Dashboard pages shall load within \textbf{3 seconds} on a 4G mobile connection (10 Mbps).
    \item The system shall support at least \textbf{100 concurrent users} without degradation in response time beyond 20\%.
    \item Database queries for referral tables shall return results in under \textbf{1 second} for datasets up to 10,000 records.
\end{itemize}

\subsection{Safety Requirements}
\begin{itemize}
    \item The system shall display a clear disclaimer on all prediction result screens stating: \textit{``This prediction is an AI-assisted triage tool and does not constitute a medical diagnosis. Clinical judgment must be applied.''}
    \item The system shall not allow a referral to be marked \texttt{completed} without a reviewer ID being recorded, ensuring accountability.
    \item Patient records must never be permanently deleted via the UI; soft-deletion or archiving must be used if removal is required.
\end{itemize}

\subsection{Security Requirements}
\begin{itemize}
    \item All data in transit must be encrypted using TLS 1.2 or higher (HTTPS enforced).
    \item Row-Level Security (RLS) must be enabled on all Supabase tables at all times.
    \item User passwords must be hashed by Supabase Auth (bcrypt); plaintext passwords must never be stored or logged.
    \item Session tokens must expire after \textbf{24 hours} of inactivity.
    \item The API must reject requests without a valid authentication token with HTTP 401.
    \item Environment variables (\texttt{.env}) containing API keys or database credentials must never be committed to version control.
\end{itemize}

\subsection{Software Quality Attributes}
\begin{itemize}
    \item \textbf{Usability:} The ANM interface must be operable by a user with minimal digital literacy. All critical actions must be completable within 3 taps. Error messages must use plain language.
    \item \textbf{Reliability:} System uptime must be 99\% during operating hours (6 AM--10 PM IST). On backend failure, the frontend must show an informative error state.
    \item \textbf{Maintainability:} The ML model must be replaceable by swapping \texttt{xgboost\_model.json} with no code changes. Database migrations must be version-controlled and applied sequentially.
    \item \textbf{Portability:} The frontend must run on any modern browser across Android, iOS, Windows, and macOS without modification.
    \item \textbf{Scalability:} The backend must support horizontal scaling via Docker containerisation.
    \item \textbf{Accessibility:} The UI must use sufficient colour contrast (WCAG AA) and support keyboard navigation on desktop.
\end{itemize}

\subsection{Business Rules}
\begin{itemize}
    \item An ANM may only create referrals to PHCs; they cannot directly refer to CHCs.
    \item Only PHC doctors may escalate referrals to CHCs.
    \item Each user can only view and manage records belonging to their facility (enforced via RLS policies).
    \item A prediction result is advisory only and must be reviewed by a qualified healthcare professional before any clinical action is taken.
    \item Patient records are owned by the facility that created them; cross-facility access is granted only through the referral mechanism.
\end{itemize}

% ─────────────────────────────────────────────
%  6. OTHER REQUIREMENTS
% ─────────────────────────────────────────────
\section{Other Requirements}

\subsection{Database Requirements}
The system uses four primary tables in Supabase PostgreSQL: \texttt{profiles}, \texttt{patients}, \texttt{predictions}, and \texttt{referrals}. All tables must have RLS policies applied. The \texttt{referrals} table implements a state machine: \texttt{pending} $\rightarrow$ \texttt{reviewed} $\rightarrow$ \texttt{in\_progress} $\rightarrow$ \texttt{completed} (or \texttt{rejected} from \texttt{pending}). Vitals are stored as JSONB for schema flexibility. All migrations are version-controlled under \texttt{/supabase/migrations/}.

\subsection{Internationalisation Requirements}
The UI must support English and Hindi as primary languages in Version 1.0. The internationalisation framework must be structured to allow addition of Marathi, Tamil, and Telugu in a future release without architectural changes. Date and time must be displayed in IST (UTC+5:30).

\subsection{Legal and Compliance Requirements}
Patient health data handling must comply with applicable Indian data protection legislation. No patient data shall be transmitted to any third-party service not explicitly listed in this SRS. The AI disclaimer (Section 5.2) is a mandatory UI element.

% ─────────────────────────────────────────────
%  APPENDIX A: GLOSSARY
% ─────────────────────────────────────────────
\newpage
\section*{Appendix A: Glossary}
\addcontentsline{toc}{section}{Appendix A: Glossary}

\begin{tabular}{|p{3cm}|p{10cm}|}
    \hline
    \textbf{Term} & \textbf{Definition} \\
    \hline
    ANM & Auxiliary Nurse Midwife -- frontline health worker at Sub-Centre level \\
    \hline
    PHC & Primary Health Centre \\
    \hline
    CHC & Community Health Centre \\
    \hline
    XGBoost & Extreme Gradient Boosting -- ML algorithm used for disease prediction \\
    \hline
    RLS & Row-Level Security -- PostgreSQL/Supabase feature restricting data access per row \\
    \hline
    SRS & Software Requirements Specification \\
    \hline
    JWT & JSON Web Token -- used for session authentication \\
    \hline
    CORS & Cross-Origin Resource Sharing \\
    \hline
    JSONB & Binary JSON storage type in PostgreSQL \\
    \hline
    UUID & Universally Unique Identifier \\
    \hline
    TBD & To Be Determined \\
    \hline
    IST & Indian Standard Time (UTC+5:30) \\
    \hline
\end{tabular}

% ─────────────────────────────────────────────
%  APPENDIX B: ANALYSIS MODELS
% ─────────────────────────────────────────────
\section*{Appendix B: Analysis Models}
\addcontentsline{toc}{section}{Appendix B: Analysis Models}

% ── TikZ styles ──────────────────────────────
\tikzset{
  state/.style={rectangle, rounded corners=6pt, draw=black, thick,
                minimum width=2.2cm, minimum height=0.8cm,
                text centered, font=\small\ttfamily, fill=gray!10},
  arrow/.style={-{Stealth[length=6pt]}, thick},
  dbbox/.style={rectangle, draw=black, thick, minimum width=3cm,
                minimum height=1cm, text centered, font=\small\bfseries,
                fill=gray!12},
  tierbox/.style={rectangle, rounded corners=4pt, draw=black, thick,
                  minimum width=3.5cm, minimum height=1.2cm,
                  text centered, font=\normalsize\bfseries, fill=gray!10},
  apibox/.style={rectangle, draw=black, thick, minimum width=2.8cm,
                 minimum height=0.9cm, text centered, font=\small},
}

% ─────────────────────────────────────────────
\subsection*{B.1 \ Referral State Machine}
% ─────────────────────────────────────────────

\begin{figure}[H]
\centering
\begin{tikzpicture}[node distance=2.6cm]

  \node[state] (pending)    {pending};
  \node[state, right=of pending]  (reviewed)   {reviewed};
  \node[state, right=of reviewed] (inprog)     {in\_progress};
  \node[state, right=of inprog]   (completed)  {completed};
  \node[state, below=1.4cm of pending] (rejected) {rejected};

  \draw[arrow] (pending)   -- node[above, font=\scriptsize]{doctor reviews}   (reviewed);
  \draw[arrow] (reviewed)  -- node[above, font=\scriptsize]{treatment starts} (inprog);
  \draw[arrow] (inprog)    -- node[above, font=\scriptsize]{case closed}       (completed);
  \draw[arrow] (pending)   -- node[right, font=\scriptsize]{rejected}          (rejected);

\end{tikzpicture}
\caption*{Referral lifecycle state machine}
\end{figure}

% ─────────────────────────────────────────────
\subsection*{B.2 \ Three-Tier System Architecture}
% ─────────────────────────────────────────────

\begin{figure}[H]
\centering
\begin{tikzpicture}[node distance=0.9cm and 1.8cm]

  % Tier boxes
  \node[tierbox, fill=blue!10]  (frontend) {Frontend (React/Vite)\\{\small Vercel}};
  \node[tierbox, fill=orange!15, right=2.2cm of frontend] (backend)  {Backend (FastAPI)\\{\small HuggingFace Spaces}};
  \node[tierbox, fill=green!12, below=1.6cm of backend]   (db)       {Database (Supabase)\\{\small PostgreSQL + RLS}};
  \node[tierbox, fill=purple!10, below=1.6cm of frontend] (browser)  {Browser\\{\small Web Speech API}};

  % Arrows
  \draw[arrow] (frontend) -- node[above, font=\scriptsize]{REST/JSON (HTTPS)} (backend);
  \draw[arrow] (frontend) -- node[left,  font=\scriptsize]{Supabase JS client} (db);
  \draw[arrow] (backend)  -- node[right, font=\scriptsize]{PostgREST} (db);
  \draw[arrow] (browser)  -- node[below, font=\scriptsize]{voice input} (frontend);

\end{tikzpicture}
\caption*{High-level system architecture}
\end{figure}

% ─────────────────────────────────────────────
\subsection*{B.3 \ Three-Tier Healthcare Data Flow}
% ─────────────────────────────────────────────

\begin{figure}[H]
\centering
\begin{tikzpicture}[node distance=3cm]

  \node[tierbox, fill=blue!10]   (anm) {ANM\\{\small Sub-Centre}};
  \node[tierbox, fill=orange!10, right=of anm] (phc) {PHC Doctor\\{\small Primary Health Centre}};
  \node[tierbox, fill=red!10,    right=of phc] (chc) {CHC Specialist\\{\small Community Health Centre}};

  \draw[arrow] (anm) -- node[above, font=\scriptsize]{refer patient} node[below, font=\scriptsize]{(low/medium/high)} (phc);
  \draw[arrow] (phc) -- node[above, font=\scriptsize]{escalate case} node[below, font=\scriptsize]{(high/critical)} (chc);

  % Feedback arrows
  \draw[arrow, dashed, bend left=30] (phc) to node[above, font=\scriptsize]{status update} (anm);
  \draw[arrow, dashed, bend left=30] (chc) to node[above, font=\scriptsize]{outcome} (phc);

\end{tikzpicture}
\caption*{Patient referral data flow across healthcare tiers}
\end{figure}

% ─────────────────────────────────────────────
\subsection*{B.4 \ Entity-Relationship Diagram}
% ─────────────────────────────────────────────

\begin{figure}[H]
\centering
\begin{tikzpicture}[node distance=2.2cm and 3.5cm]

  \node[dbbox, fill=blue!10]   (profiles)    {profiles\\{\scriptsize id, full\_name, role,}\\{\scriptsize facility\_name}};
  \node[dbbox, fill=green!10,  right=3.5cm of profiles] (patients)    {patients\\{\scriptsize id, patient\_name, age,}\\{\scriptsize gender, symptoms}};
  \node[dbbox, fill=orange!10, below=2cm of patients]   (predictions) {predictions\\{\scriptsize id, predicted\_disease,}\\{\scriptsize confidence}};
  \node[dbbox, fill=red!10,    below=2cm of profiles]   (referrals)   {referrals\\{\scriptsize id, priority, status,}\\{\scriptsize referred\_to}};

  % Relationships
  \draw[arrow] (profiles)    -- node[above, font=\scriptsize]{creates}   (patients);
  \draw[arrow] (patients)    -- node[right, font=\scriptsize]{has}        (predictions);
  \draw[arrow] (patients)    -- node[above, font=\scriptsize]{triggers}   (referrals);
  \draw[arrow] (profiles)    -- node[left,  font=\scriptsize]{refers}     (referrals);
  \draw[arrow] (predictions) -- node[below, font=\scriptsize]{linked to}  (referrals);

\end{tikzpicture}
\caption*{Entity-relationship diagram (simplified)}
\end{figure}

% ─────────────────────────────────────────────
\subsection*{B.5 \ API Endpoint Summary}
% ─────────────────────────────────────────────

\begin{center}
\begin{tabular}{|l|l|l|l|}
    \hline
    \textbf{Method} & \textbf{Endpoint} & \textbf{Request} & \textbf{Response} \\
    \hline
    GET  & \texttt{/health}   & --                      & \texttt{\{status, model\_loaded\}} \\
    \hline
    GET  & \texttt{/symptoms} & --                      & \texttt{\{symptoms[], count\}} \\
    \hline
    POST & \texttt{/predict}  & \texttt{\{symptoms[]\}} & \texttt{\{disease, confidence, description, precautions[]\}} \\
    \hline
\end{tabular}
\end{center}

% ─────────────────────────────────────────────
%  APPENDIX C: TO BE DETERMINED LIST
% ─────────────────────────────────────────────
\section*{Appendix C: To Be Determined List}
\addcontentsline{toc}{section}{Appendix C: To Be Determined List}

\begin{enumerate}
    \item \textbf{TBD-01:} Specific regional language (Marathi, Tamil, Telugu) support -- implementation timeline not yet defined.
    \item \textbf{TBD-02:} SMS/WhatsApp notification on referral status change -- third-party service selection pending.
    \item \textbf{TBD-03:} Offline mode capability for ANMs in areas with no internet connectivity -- feasibility study not yet conducted.
    \item \textbf{TBD-04:} Integration with government health portals (e.g., ABDM/ABHA) -- regulatory requirements under review.
    \item \textbf{TBD-05:} Model retraining pipeline and dataset update cadence -- data governance policy not yet established.
\end{enumerate}

\end{document}
